\begin{textsample}{POS Dim 3 – 2023\_c6 – Score 2.00 – 2023\_c6/enade\_2023\_nutri\_ab5bde7002.txt}  \label{ex:f3_pos_109}
CAPÍTULO 5 ANÁLISE TÉCNICA DA PROVA Este capítulo tem o objetivo de apresentar o desempenho dos estudantes concluintes de Nutrição no Enade 2023.
Para isso, foram calculadas as estatísticas básicas da prova em seu todo ( subseção 5.1.1 ), bem como as estatísticas relacionadas aos componentes de Formação Geral ( subseção 5.1.2 ) e de Conhecimento Específico da área ( subseção 5.1.3 ).
Dadas as suas características, foram analisadas, em separado, as questões objetivas ( seção 5.2 ) e as questões discursivas ( seção 5.3 ).
Tomando -se como base a questão discursiva do componente de Formação Geral, nas subseções de 5.3.1.4 a 5.3.1.6, são apresentados os resultados e os comentários sobre a correção das respostas em relação ao desempenho linguístico.
Para as questões objetivas, foram disponibilizados os índices de facilidade e de discriminação ponto-bisserial, também separadamente, para os componentes de Formação Geral ( subseção 5.2.1 ) e de Conhecimento Específico ( subseção 5.2.2 ).
De cada componente, uma das questões foi escolhida para exemplificar a análise gráfica, relacionando -se as alternativas escolhidas pelos estudantes ( inclusive o gabarito ) ao número de acertos no componente.
No Anexo IV, apresenta -se a íntegra da análise gráfica para todas as questões objetivas.
Para cada uma das questões discursivas, os conteúdos dos tipos mais comuns de respostas dos estudantes são apresentados e comparados com o padrão de resposta esperado ( ver Anexo VIII ).
Nas tabelas, constam as seguintes estatísticas das notas⁵ : média do desempenho na prova, erro-padrão da média, desvio-padrão, nota mínima, mediana e nota máxima para cada um de seus componentes.
Tais estatísticas contemplam o total de estudantes concluintes da área de Nutrição inscritos e presentes na prova do Enade 2023, tendo em vista agregações, ou por grandes regiões e o país como um todo, ou por categoria administrativa, organização acadêmica da IES e modalidade de oferta dos cursos.
Em relação aos gráficos de distribuição de notas, o intervalo considerado foi de 10 unidades, aberto à esquerda e fechado à direita, com exceção do primeiro intervalo, [ 0 ; 10 ], fechado em ambos os extremos.
Para os gráficos de distribuição das notas das questões discursivas, foram consideradas mais duas categorias : questão em branco⁶ e nota zero. 5 Uma definição dessas estatísticas pode ser encontrada no Glossário inserido no final desse relatório. 6 Nesse grupo estão incluídas também as respostas classificadas como nulas ou desconsideradas. 5.1 ESTATÍSTICAS BÁSICAS DA PROVA Esta seção apresenta estatísticas selecionadas e histogramas da nota geral ( 5.1.1 ) e de cada componente : Formação Geral ( 5.1.2 ) e Conhecimento Específico ( 5.1.3 ).
São também apresentadas estatísticas selecionadas de subpopulações, caracterizadas por grande região, categoria administrativa, organização acadêmica e modalidade de oferta. 5.1.1 ESTATÍSTICAS BÁSICAS GERAIS Na Tabela 5.1, são apresentadas as estatísticas básicas da prova, por grande região, dos estudantes concluintes de Nutrição.
Como se pode verificar, a partir dos dados mostrados na Tabela 5.1, a média das notas da prova foi de 48,7 ( nas seções seguintes serão analisados os componentes de Formação Geral e de Conhecimento Específico ).
Os estudantes da região Norte obtiveram a média mais baixa ( 44,9 ), e os da região Centro-Oeste alcançaram a média mais alta ( 49,7 ).
As demais médias foram 48,3 na região Nordeste ; 49,5 na região Sudeste ; e 48,1 na região Sul.
O desvio-padrão para o Brasil foi de 13,3, sendo o maior encontrado na região Nordeste ( 13,5 ), e o menor, na região Norte ( 12,9 ).
A região que obteve a maior nota máxima foi a região Sudeste ( 91,3 ), ao passo que a região que atingiu a menor nota máxima foi o Norte ( 75,7 ).
A mediana do Brasil, como um todo, foi igual a 49,1, sendo a maior mediana obtida na região Sudeste ( 50,1 ), e a menor, na região Norte ( 45,0 ).
Considerando -se as notas segundo grande região, observa -se que existe diferença estatisticamente significativa ao nível de 95 \% entre as médias obtidas em todas as regiões, exceto entre Sudeste e Centro-Oeste ; Sul e Nordeste.
O comportamento das notas dos estudantes de todo o Brasil pode ser observado no Gráfico 5.1.
O Gráfico 5.1 apresenta um histograma com a distribuição das notas dos estudantes de Nutrição de todo o Brasil : uma distribuição unimodal com a moda no intervalo ( 50 ; 60 ].
Na Tabela 5.2, são apresentadas informações referentes à média da nota final, desagregadas por categoria administrativa, organização acadêmica e modalidade de oferta dos cursos.
Os dados da Tabela 5.2 mostram que, em relação à categoria administrativa, os estudantes das IES pública obtiveram média mais alta ( 58,2 ) que a média das IES privada ( 47,1 ).
O desvio-padrão para as IES públicas foi de 11,6, indicando uma dispersão baixa das notas nesta categoria administrativa.
Já o desvio-padrão da IES privadas foi 12,9.
No que se refere à organização acadêmica, os estudantes de universidades obtiveram a média de 50,4, enquanto os de centros universitários, faculdades e CEFET/IF obtiveram, respectivamente, médias de 47,6, 44,8 e 51,8.
Constata -se que existe diferença estatisticamente significativa, ao nível de 95 \%, entre as médias das organizações acadêmicas, exceto entre universidade e CEFET/IF.
A média da modalidade a distância ( 45,8 ) foi inferior à média da modalidade presencial ( 49,6 ), sendo a diferença entre as duas estatisticamente significativa, ao nível de 95 \%. 5.1.2 ESTATÍSTICAS BÁSICAS NO COMPONENTE DE FORMAÇÃO GERAL Na Tabela 5.3, são apresentadas as estatísticas básicas em relação ao componente da prova que avalia a Formação Geral dos estudantes concluintes.
A nota de Formação Geral é considerando -se os acertos nas nove questões objetivas e a nota média obtida na questão discursiva, sendo que a nota dessa questão obtida leva em conta 80 \% da nota relativa ao conteúdo e 20 \% da nota relativa ao desempenho linguístico.
Conforme os dados da Tabela 5.3, os estudantes de todo o Brasil obtiveram, no componente de Formação Geral da prova, desempenho médio de 48,4.
Quanto à variabilidade, o desvio-padrão das notas dos estudantes do Brasil foi 16,6.
A maior média foi obtida na região Sudeste ( 49,2 ), e a menor, na região Norte ( 46,2 ).
As demais médias foram : 48,8 na região Centro-Oeste ; 48,6 na região Nordeste ; e 47,1 na região Sul.
Já o maior desvio-padrão foi obtido na região Nordeste ( 16,8 ), e o menor, nas regiões Centro-Oeste, Norte e Sudeste ( 16,5 ).
Os demais desvios foram : 16,7 na região Sul.
A maior nota no componente de Formação Geral da prova do Enade 2023 foi 100,0, obtida por pelo menos um estudante na região Sudeste.
A menor nota máxima foi obtida na região Norte ( 90,9 ).
A mediana do Brasil, como um todo, foi de 48,5 ; sendo a menor mediana encontrada na região Norte ( 46,8 ), e a maior nas regiões Centro-Oeste e Sudeste ( 49,3 ).
A nota mínima nessa parte foi zero em quase todas as regiões, com exceção da região Centro-Oeste, em que a nota foi 5,4. e a maior nas regiões Centro-Oeste e Sudeste ( 49,3 ).
A nota mínima nessa parte foi zero em quase todas as regiões, com exceção da região Centro-Oeste, em que a nota foi 5,4.
Considerando -se as notas, segundo grande região, observa -se que existe diferença estatisticamente significativa entre as médias das notas do componente de Formação Geral obtidas nas regiões, exceto entre as regiões Nordeste e Centro-Oeste ; Sudeste e Centro-Oeste ; Sudeste e Nordeste ; Sul e Norte.
No Gráfico 5.2, é possível verificar a avaliação do desempenho dos estudantes no componente de Formação Geral, a partir do histograma da distribuição das notas correspondentes.
Nesse gráfico, a distribuição é unimodal, com moda em ( 40 ; 50 ], mas com pouca diferença para o intervalo que o precede ( 50:60 ].
Na Tabela 5.4, são apresentadas as informações referentes ao desempenho dos concluintes no componente de Formação Geral, em diferentes agregações : categoria administrativa, organização acadêmica e modalidade de oferta.
A partir dos dados da Tabela 5.4, observa -se que há diferença estatisticamente significativa entre as médias dos tipos de categoria administrativa.
A maior média foi obtida por estudantes de IES pública ( 58,4 ), e a menor, pelos de IES privada ( 46,7 ).
Em relação às médias por tipos de organização acadêmica, encontra -se : 50,1 para as universidades ; 47,3 para os centros universitários, 44,5 para as faculdades e 53,6 para os CEFET/IF.
A nota máxima obtida pelas IES privadas ( 100,0 ) foi maior em comparação às IES públicas.
O desvio-padrão menor foi percebido nas IES públicas ; e, entre as organizações acadêmicas, esse valor foi menor para os CEFET/IF.
Já no que se refere à modalidade de oferta, os cursos presenciais apresentaram média de 49,3, enquanto a média dos cursos a distância foi de 45,6.
A diferença é, portanto, estatisticamente significativa. 5.1.3 ESTATÍSTICAS BÁSICAS DO COMPONENTE DE CONHECIMENTO ESPECÍFICO Na Tabela 5.5, são apresentadas as estatísticas básicas referentes ao componente de Conhecimento Específico da área de Nutrição.
A nota do componente de Conhecimento Específico leva em conta as notas da parte objetiva e a média das notas obtidas na questão da parte discursiva da área.
De acordo com a Tabela 5.5, a média do desempenho dos estudantes do Brasil, para o componente de Conhecimento Específico da prova, foi de 48,8.
A maior média foi obtida na região Centro-Oeste ( 49,9 ), e a menor, na região Norte ( 44,5 ).
As demais médias foram : 48,1, na região Nordeste ; 49,6, na região Sudeste ; e 48,4, na região Sul.
Quanto à variabilidade das notas, o desvio-padrão do Brasil foi 14,4 ; sendo o maior desvio-padrão observado nas regiões Centro-Oeste e Nordeste ( 14,5 ), e o menor, na região Norte ( 14,1 ).
Os demais desvios foram : 14,3 na região Sudeste ; e 14,2 na região Sul.
A mediana das notas dos estudantes de todo o Brasil foi de 49,0.
A maior mediana ocorreu na região Centro-Oeste ( 50,9 ), e a menor, na região Norte ( 44,2 ).
As demais medianas foram : 48,0 na região Nordeste ; 50,5 na região Sudeste ; e 48,5 na região Sul.
A nota máxima do Brasil, como um todo, foi 96,1 ; sendo obtida por, pelo menos, um estudante na região Nordeste.
As demais notas máximas foram : 88,3 na região Centro-Oeste ; 80,5 na região Norte ; 95,1 na região Sudeste ; e 89,2 na região Sul.
A nota mínima nessa parte foi zero nas regiões Nordeste e Sudeste.
Nas regiões Centro-Oeste, Norte e Sul, as notas mínimas foram, respectivamente, 7,4 ; 3,9 ; e 3,9.
Observa -se que existe diferença estatisticamente significativa entre todas as regiões na análise das médias das notas do componente de Conhecimento Específico, exceto entre as regiões Sudeste e Centro-Oeste ; Sul e Nordeste.
Assim como nos Gráficos 5.1 e 5.2, mostrados anteriormente, no Gráfico 5.3, apresentado acima, foi possível ser realizada uma avaliação do desempenho de concluintes de Nutrição, em relação ao componente de Conhecimento Específico, com um histograma da distribuição das notas correspondentes.
Esta também é uma distribuição unimodal, e o grupo modal é o ( 50 ; 60 ].
Na Tabela 5.6, é apresentada uma comparação dos resultados em relação à categoria administrativa, à organização acadêmica e à modalidade de oferta, levando -se em conta o desempenho de estudantes do componente de Conhecimento Específico da prova.
De acordo com os dados expostos na Tabela 5.6, no que se refere à organização acadêmica, a maior média foi dos CEFET/IF ( 51,2 ), seguida das universidades ( 50,5 ), dos centros universitários ( 47,6 ) e, por fim, das faculdades ( 44,9 ).
O maior desvio-padrão foi o das universidades ( 14,7 ).
Os CEFET/IF obtiveram a maior mediana ( 52,0 ).
Já as notas máximas foram : 95,1 nas universidades, 96,1 nos centros universitários, 85,7 nas faculdades e 76,8 nos CEFET/IF.
A nota mínima foi zero para as universidades e para os centros universitários, sendo 3,9 para faculdades e 24,0 para CEFET/IF.
Observa -se que existe diferença estatisticamente significativa, ao nível de 95 \%, no componente de Conhecimento Específico entre as notas dos tipos de organização acadêmica, exceto entre universidade e CEFET/IF.
Quanto à categoria administrativa, existe diferença estatisticamente significativa entre as médias das IES públicas ( 58,1 ) e as das IES privadas ( 47,2 ).
Neste caso, assim como no componente de Formação Geral, a maior média foi obtida por estudantes de IES públicas de ensino.
Os cursos de educação presencial apresentaram maior média ( 49,7 ) em relação aos cursos a distância ( 45,8 ), sendo a diferença estatisticamente significativa. 5.2 ANÁLISE DAS QUESTÕES OBJETIVAS Esta seção apresenta estatísticas selecionadas e histogramas das questões objetivas dos componentes de Formação Geral ( 5.2.1 ) e de Conhecimento Específico ( 5.2.2 ).
São, também, apresentadas e comparadas as médias das subpopulações caracterizadas por grande região. 5.2.1 COMPONENTE DE FORMAÇÃO GERAL Na Tabela 5.7, são apresentadas as estatísticas básicas relativas às nove questões objetivas do componente da prova que abrange a Formação Geral dos estudantes.
De acordo com os dados da Tabela 5.7, em relação ao componente de Formação Geral, a média do Brasil foi 51,7.
A menor média foi encontrada na região Norte ( 48,3 ), e a maior, na região Sudeste ( 52,8 ).
As demais médias foram : 52,3 na região Centro-Oeste ; 51,6 na região Nordeste ; e 50,4 na região Sul.
O desvio-padrão do Brasil foi 19,1 ; sendo o maior desvio-padrão encontrado na região Sul ( 19,3 ), e o menor desvio-padrão encontrado nas regiões Centro-Oeste, Nordeste e Norte ( 19,0 ).
A mediana do Brasil, como um todo, foi 55,6 e das regiões Sul, Nordeste, Centro-Oeste, Norte e Sudeste foram iguais a 55,6, 55,6, 55,6, 44,4 e 55,6, respectivamente.
Quanto à nota máxima, todas as regiões alcançaram a nota 100,0.
Já as notas mínimas foram 0,0, também em todas as regiões.
Na Tabela 5.8, exposta mais adiante, são apresentados o índice de facilidade e o índice de discriminação ( ponto-bisserial ) para cada uma das questões objetivas do componente de Formação Geral.
De acordo com os dados da Tabela 5.8, as questões objetivas do componente de Formação Geral, segundo o índice de facilidade, foram assim avaliadas : nenhuma teve o índice de facilidade classificado como “ Muito difícil ” ; nenhuma questão foi classificada como “ Muito fácil ” ; duas questões foram classificadas como “ Fácil ” ; cinco questões foram classificadas como"Médio ” ; e duas questões foram classificadas na categoria “ Difícil ”.
O índice de facilidade variou de 0,24 a 0,78.
Já em relação ao índice de discriminação, cinco questões foram classificadas com o índice “ Muito bom ” ; três com o índice “ Bom ”, uma questão com o índice “ Médio ” ; e nenhuma com o índice"Fraco ”.
O índice de discriminação variou de 0,27 a 0,49.
Na Tabela 5.9, é apresentada a distribuição das questões, levando -se em conta, simultaneamente, a classificação dos dois índices.
Os dados da Tabela 5.9, em conjunto com os dados relativos às questões objetivas da Tabela 5.8, atestam que cinco questões que tiveram índice de discriminação “ Muito bom ” figuraram em dois níveis de dificuldade – “ Médio e Fácil ” : nenhuma classificada na categoria “ Muito Fácil ” ; uma classificadas na categoria “ Fácil ” ( Questão 6 ) ; quatro classificadas na categoria “ Médio ” ( Questões 1, 3, 4 e 8 ) ; nenhuma, na categoria “ Difícil ” ; e nenhuma, na categoria “ Muito difícil ”.
Em particular, a Questão 6 foi a que apresentou o maior poder discriminatório, com índice 0,49, considerado “ Fácil ” em termos de facilidade, com uma proporção de 0,71 acertos, não tendo essa questão alcançado o máximo de acertos.
No Gráfico 5.4, para exemplificar, analisa -se o comportamento da questão objetiva 4 de Formação Geral.
Trata -se de uma questão considerada com nível “ Médio ” de dificuldade, em relação à facilidade, e que obteve o índice de discriminação 0,45 ( muito bom ).
No Gráfico 5.4, cada uma das cinco curvas representa o percentual de respostas em determinada alternativa da questão, em função do número de acertos dos estudantes nessa parte da prova ( Formação Geral/múltipla escolha ), antes de possíveis eliminações pelo critério ponto-bisserial.
A curva em vermelho corresponde à alternativa A, a alternativa correta para essa questão.
Observa -se que, entre os estudantes com menor número de acertos, nessa parte do exame, a situação mais frequente foi a escolha da alternativa B ( em mostarda ), incorreta.
Por exemplo, entre os estudantes que acertaram até duas respostas, 13,2 \% escolheram a alternativa A ( em vermelho ) ; 45,5 \% escolheram a alternativa B ( em mostarda ) ; 5,5 \% escolheram a alternativa C ( em verde ) ; 11,4 \% escolheram a alternativa D ( em azul ) ; e 23,9 \% escolheram a alternativa E ( em lilás ).
Entre os que acertaram duas respostas entre as questões de múltipla escolha de Formação Geral, 0,4 \% deixou a questão em branco e 0,1 \% marcou mais de uma alternativa, invalidando a questão.
À medida que o número de acertos aumenta, indicando desempenho melhor nessa parte da prova, aumenta, concomitantemente, a proporção de estudantes que selecionaram a alternativa correta A, atingindo 100 \% para os estudantes com oito acertos.
Essa 5.2.2 COMPONENTE DE CONHECIMENTO ESPECÍFICO Na Tabela 5.10, são apresentadas as estatísticas básicas em relação às questões objetivas do componente de Conhecimento Específico da prova de Nutrição por grande região.
Em relação ao componente de Conhecimento Específico da prova de Nutrição, como será discutido ainda nesta seção, seis questões não foram usadas no cômputo das notas por terem sido descartadas pelo critério do ponto-bisserial.
Assim, as notas foram calculadas com base em 23 das 29 questões objetivas válidas do componente de Conhecimento Específico.
Como se verifica a partir dos dados da Tabela 5.10, a média do Brasil para o componente de Conhecimento Específico da prova foi 51,1.
A menor média foi observada na região Norte ( 46,3 ), e a maior, na região Centro-Oeste ( 52,2 ).
O desvio-padrão de todo o Brasil foi 14,9, sendo o maior desvio-padrão encontrado na região Nordeste ( 15,1 ), e o menor, na região Norte ( 14,6 ).
A mediana de todo o Brasil foi 52,2, o mesmo valor da mediana encontrada em todas as regiões, com exceção da região Norte, que obteve mediana 47,8.
A nota máxima da prova foi 95,7, obtida, nas questões objetivas do componente de Conhecimento Específico, por, pelo menos, um estudante das regiões Nordeste e Sudeste.
Nas demais regiões, a nota máxima variou entre 82,6 ( Norte ) e 91,3 ( Sul e Centro-Oeste ).
A nota mínima de maior valor foi 4,3, obtida por, pelo menos, um estudante nas regiões Centro-Oeste, Norte e Sul.
Na Tabela 5.11, são apresentados os índices de facilidade e de discriminação ( ponto-bisserial ) das questões objetivas do componente de Conhecimento Específico, para os estudantes de Nutrição.
A partir do índice de facilidade obtido, apresentado na Tabela 5.11, pode -se concluir que a maioria das questões objetivas da prova foram classificadas no índice “ Fácil ” ( 12 ). 11 questões foram classificadas no índice"Difícil ” ; cinco, no índice “ Médio ” ; nenhuma, no índice “ Muito fácil ” ; e uma questão foi alocada no índice “ Muito difícil ”.
Quanto aos índices de discriminação das questões objetivas do componente de Conhecimento Específico da prova, obtém -se como resultado a seguinte classificação : duas questões foram alocadas no índice “ Muito bom ” e nove, no “ Bom ”.
Dessa forma, 11 das 29 questões foram classificadas com os índices de discriminação “ Bom ” e “ Muito bom ”.
Entre as demais, dez delas foram alocadas no índice “ Médio ” e oito, no “ Fraco ”.
Constata -se, assim, que a prova – no que se refere ao componente de Conhecimento Específico – possuía média capacidade de discriminar entre estudantes que dominam ou não o conteúdo.
O índice de facilidade variou de 0,14 a 0,81, e o de discriminação, de 0,02 a 0,44.
Na Tabela 5.12, é apresentada a distribuição das questões considerando -se, simultaneamente, a classificação dos dois índices.
Os dados da Tabela 5.12 atestam que nove questões que tiveram índice de discriminação “ Bom ” figuraram em três níveis de dificuldade – “ Difícil","Fácil"e"Médio ” : seis classificadas na categoria “ Fácil ” ; duas na categoria “ Difícil ” ; uma na categoria “ Médio ” ; nenhuma nas categorias “ Muito fácil ” e “ Muito difícil ”.
As questões que tiveram o índice de discriminação “ Fraco ” configuraram -se com índices de facilidade “ Muito fácil ” ( nenhuma questão ), “ Fácil ” ( uma questão ), “ Médio ” ( nenhuma questão ), “ Difícil ” ( seis questões ) e “ Muito Difícil ” ( uma questão ).
Entre as questões que alcançaram os maiores índices de discriminação, “ Muito bom ”, duas delas, as Questões 32 e 34, tiveram o maior índice de discriminação ( 0,44 ).
Tais questões foram classificadas como “ Fácil ”, com índice de facilidade igual a 0,62 e 0,63, respectivamente.
A Questão 11 foi a mais “ Difícil ” entre as 29 questões específicas válidas, com baixo índice de facilidade, apenas 0,14 de proporção de acertos.
Essa questão apresentou poder discriminatório fraco, 0,19, o que comprova ter sido a mais “ Difícil ” para os estudantes.
Destacam -se, também, com baixo índice de facilidade, as Questões 22, 30 e 11, cujos índices de facilidade, em termos proporcionais, correspondem, respectivamente, a 0,20, 0,19 e 0,14 de estudantes que responderam acertadamente.
Essas três questões também tiveram poder discriminatório “ Fraco","Fraco"e"Fraco ”, com índices 0,09, 0,02 e 0,19, respectivamente.
Destaca -se que as Questões 11, 13, 15, 22, 23 e 30 foram consideradas inadequadas pelo critério ponto-bisserial.
Por isso, foram eliminadas do cômputo da nota final.
A título de exemplo das análises do comportamento das questões objetivas, no Gráfico 5.5, analisa -se a Questão 32 do componente de Conhecimento Específico.
No Gráfico 5.5, cada uma das cinco curvas representa o percentual de respostas em determinada alternativa da questão, em função do número de acertos dos estudantes nessa parte da prova ( componente de Conhecimento Específico ), antes de possíveis eliminações pelo critério ponto-bisserial.
A alternativa correta D, representada no gráfico pela curva em azul, foi escolhida em maiores proporções pelos estudantes com desempenho melhor nessa parte da prova.
Já as alternativas incorretas, também denominadas distratores, foram selecionadas, principalmente, por aqueles com notas mais baixas.
A Questão 32 foi considerada, pelas respostas dos estudantes avaliados na prova, como uma questão com índice de facilidade “ Fácil, 0,62 ”, ou seja, 62,0 \% dos estudantes assinalaram, acertadamente, a opção D, correspondente ao gabarito.
Seu índice de discriminação foi igual a 0,44, classificado como “ Muito bom ”, o maior valor de discriminação.
Observa -se que a soma não é 100 \%, em razão das questões não respondidas ou com mais de uma opção marcada.
Cerca de 83,3 \% daqueles com nota zero marcaram a opção B ( linha mostarda ).
A proporção de estudantes que selecionaram a resposta correta D aumenta gradativamente, chegando a atingir 100 \% para 22 acertos, enquanto a proporção dos que escolheram alternativas incorretas decai, a partir de sete acertos, em função do número de acertos nessa parte da prova.
Os gráficos relativos às demais questões do componente de Conhecimento Específico constam no Anexo IV. 5.3 ANÁLISE DAS QUESTÕES DISCURSIVAS Esta seção apresenta estatísticas selecionadas e histogramas das questões discursivas de Formação Geral ( 5.3.1 ) e de Conhecimento Específico ( 5.3.2 ).
Expõe e compara, também, as médias de subpopulações, caracterizadas por grande região.
Cumpre notar que nem todas as questões passam pelo crivo dos corretores.
Após a digitalização, as respostas passam por um pré-processamento de identificação automático.
Algumas dessas questões são definidas como um “ branco automático ” ( que não são enviadas para os professores para correção, a não ser que, na segunda fase, a de inspeção visual, seja constatada uma resposta ).
O critério usado nesse pré-processamento consiste em verificar, para cada questão, a partir da quantidade de pixels, o que, provavelmente, é uma resposta em branco.
A partir da separação, os documentos que foram considerados em branco são enviados para uma verificação visual, feita por uma equipe de colaboradores devidamente capacitados.
Desse modo, qualquer erro, presente no pré-processamento, é encontrado, e as questões preenchidas são enviadas para a correção dos avaliadores.
Eventualmente, algumas respostas em branco não são detectadas nesse pré-processamento automático.
A presença de pequenos riscos, ou manchas, pode fazer o programa não identificar uma questão como um possível branco.
Essas são enviadas para a correção, mas são classificadas como “ branco ” pelos avaliadores, no momento da correção. 5.3.1 COMPONENTE DE FORMAÇÃO GERAL As análises dos resultados de desempenho dos estudantes de Nutrição, na questão discursiva relativa à Formação Geral, encontram -se na Tabela 5.13 e no Gráfico 5.6, cujos dados são analisados de forma sequenciada.
No caso da Tabela 5.13, analisa -se a média das notas da questão discursiva de Formação Geral, que leva em conta as correções realizadas por duas bancas : uma que analisa e pontua aspectos relativos ao conteúdo das respostas, seguindo o padrão de respostas divulgado pelo Inep ( ver Anexo VIII ) ; e outra que avalia o desempenho linguístico expresso pela redação das respostas, seguindo um padrão de respostas específico da área de Língua Portuguesa, detalhado na seção 5.3.1.5 deste relatório.
No cálculo da nota da questão discursiva, consideram -se 80 \% relativos à nota da correção de conteúdo e 20 \% referentes à nota de desempenho linguístico.
Na Tabela 5.13, observa -se que os estudantes de todo o Brasil obtiveram, em Formação Geral, média de 42,4 nas questões discursivas.
Pode -se notar, também, que o desvio-padrão nesse conjunto de questões foi de 24,6.
A maior média foi obtida na região Nordeste ( 43,1 ), e a menor na região Sul ( 41,2 ).
A mediana de todo o Brasil, neste componente, foi 42,2 e na região Norte ( 43,5 ) foi encontrada a maior mediana.
A nota máxima ( 100,0 ) foi obtida por, pelo menos, um estudante das regiões Centro-Oeste, Nordeste, Sudeste e Sul.
A nota mínima ( 0,0 ) foi a mesma em todas as regiões do Brasil.
No Gráfico 5.6, está representada a distribuição das notas na questão discursiva do componente de Formação Geral.
Como mostra o Gráfico 5.6, a moda da distribuição ocorre no intervalo ( 50 ; 60 ], com frequência de 14,2 \% dos participantes, sendo que, no intervalo [ 0 ; 10 ], inclui -se, além da nota zero, a frequência de estudantes que deixaram esse tipo de questão em branco.
Na sequência, os resultados verificados para a questão discursiva de Formação Geral serão apresentados, estabelecendo -se relações com os conteúdos abordados.
Os comentários da banca de docentes avaliadores, a respeito do observado na correção das respostas dos estudantes, suas impressões e conclusões serão apresentadas junto à análise da questão.
Cumpre esclarecer que, tendo em vista que a questão discursiva de Formação Geral foi padronizada, ou seja, foi a mesma em todas as provas, os comentários da banca são os mesmos para todas as carreiras acadêmicas, sendo direcionados a todos os estudantes que participaram do Enade 2023.
A seguir, serão analisados o desempenho linguístico dos estudantes da área de Nutrição na questão discursiva de Formação Geral do Enade 2023, apresentando -se os resultados obtidos e os comentários da banca de correção sobre a questão. 5.3.1.1 ANÁLISE DE CONTEÚDO DA QUESTÃO DISCURSIVA DO COMPONENTE DE FORMAÇÃO GERAL Os dados de Nutrição, obtidos a partir das respostas à questão discursiva do componente de Formação Geral, encontram -se na Tabela 5.14 e no Gráfico 5.7, exposto após a apresentação dos dados da referida tabela.
Verifica -se, a partir dos dados da Tabela 5.14, que, no conteúdo da questão discursiva do componente de Formação Geral, os estudantes de todo o Brasil obtiveram média de 37,3.
A maior média foi obtida na região Nordeste ( 38,2 ), e a menor na região Sul ( 36,0 ).
Quanto à variabilidade das notas, o desvio-padrão de todo o Brasil foi 28,0.
O menor desvio-padrão foi obtido nas regiões Centro-Oeste e Sul ( 27,6 ), e o maior desvio-padrão foi obtido pela região Nordeste ( 28,6 ).
A mediana do Brasil, como um todo, foi 37,5, sendo igua à das regiões Centro-Oeste ( 37,5 ), Nordeste ( 37,5 ), Norte ( 37,5 ) e Sudeste ( 37,5 ), e superior à mediana da região Sul ( 35,0 ).
A nota máxima da questão discursiva foi 100,0, obtida por, pelo menos, um estudante em todas as regiões.
Já a nota mínima para todas as regiões foi zero.
No Gráfico 5.7, apresenta -se a distribuição das notas de conteúdo da questão discursiva do componente de Formação Geral.
O percentual de participantes que deixaram a prova em branco ou zeraram a questão discursiva na avaliação de conteúdo foi 23,4 \%.
O Gráfico 5.7 apresenta que a moda dessa distribuição ocorre no intervalo [ 0 ; 10 ], com uma frequência de 28,2 \% dos participantes. 5.3.1.2 COMENTÁRIOS SOBRE A CORREÇÃO DE CONTEÚDO DAS RESPOSTAS À QUESTÃO DISCURSIVA A questão discursiva do componente de Formação Geral foi composta por um texto motivador e por dois itens ( a e b ).
O texto motivador apresentou um comentário sobre uma publicação do Instituto Brasileiro de Geografia e Estatística ( IBGE ) que tratava dos indicadores sobre a realidade social brasileira, apontando o crescimento do nível de desigualdade socioeconômica em 2021 e concluindo que a pobreza teve aumento recorde nesse período.
A questão foi dividida em dois itens.
O primeiro item solicitava que o respondente estabelecesse uma relação entre o perfil da população mais afetada pelas desigualdades sociais e as situações de risco socioambiental.
Já o segundo item pedia a proposição de duas possíveis ações a serem desenvolvidas em bairros periféricos a fim de minimizar os impactos socioambientais.
A questão foi pertinente, apresentou linguagem acessível e abordou conteúdo basilar para Formação Geral ; no entanto, apresentou um nível de dificuldade alto, dadas as competências intrincadas que demandou.
O tema foi atual e de grande relevância para a sociedade brasileira.
Abordou assunto de impacto social e de provável domínio dos concluintes de curso superior.
O texto motivador foi adequado e coerente com o que se pediu na questão, tratou do aumento da desigualdade social e da descrição do perfil da população com base em critérios socioeconômicos.
Com relação às habilidades cognitivas, observou -se que as seguintes foram necessárias para responder à questão proposta : capacidade de expressar ideias de maneira clara, coesa e argumentativa por meio da produção de textos ; capacidade de inter-relacionar objetos de conhecimento nas diferentes áreas ; e capacidade de criar estratégias de ação, selecionar métodos e formular propostas.
Entre os saberes necessários para o desenvolvimento das respostas, os respondentes deveriam conhecer o meio social e/ou geográfico da população brasileira de perfil socioeconômico pobre ou extremamente pobre e deveriam saber relacionar esse meio com situações de risco socioambiental que afligem principalmente esse perfil populacional.
Portanto, deveriam, também, conhecer uma variedade de riscos socioambientais a que esse perfil populacional está sujeito, por conta do meio social e/ou geográfico em que vive.
A questão, em seu item b ), por sua vez, demandou o domínio de um rol de propostas minimamente factíveis para reduzir tais situações de risco, envolvendo distintos atores sociais.
Em relação ao padrão de resposta, esperava -se para cada um dos itens avaliados que : a ) O estudante explicasse a relação entre desigualdade social e riscos socioambientais percorrendo o seguinte trajeto teórico-argumentativo : • descrever o(s) risco(s) socioambiental(is) que afeta(m) a população com perfil pobre ou extremamente pobre, tais como inundações, deslizamentos de terra, contaminação ambiental, maior vulnerabilidade em relação a doenças, discriminação social e criminalidade, entre outros ; • descrever o meio ( geográfico ou social ), como morros, áreas próximas a rios, mangues, espaços marcados pela violência urbana, entre outros, nas cidades brasileiras, como fator de criação ou potencialização das condições de vulnerabilidade e risco aos pobres ou extremamente pobres ; • estabelecer nexo causal entre a pobreza ou extrema pobreza, o meio ( geográfico ou social ) urbano e a situação de risco socioambiental, relacionando esses aspectos ao perfil socioeconômico da população.
Portanto, esperava -se que o estudante apresentasse consistência argumentativa quanto ao que foi solicitado, notadamente quanto aos riscos socioambientais, à descrição do meio e ao estabelecimento do nexo causal entre pobreza, meio e risco socioambiental. b ) O estudante apresentasse propostas pertinentes, factíveis e bem desenvolvidas que envolvessem ação governamental ( federal, estadual/distrital ou municipal ) e participação da comunidade a fim de minimizar riscos socioambientais, como, por exemplo : • implantar obras de infraestrutura urbana que envolvessem contenção de morros e encostas ; saneamento básico : serviços regulares de abastecimento de água, coleta e tratamento de esgotos, coleta e manejo de resíduos sólidos, manejo de águas pluviais, entre outras ; • desenvolver políticas sociais, como segurança pública, provisão de moradia adequada, com custo acessível, e regularizada do ponto de vista fundiário ; educação de qualidade ; segurança Relatório Síntese de Área – ENADE/2023 pública ; atendimento psicossocial, com especial atenção para grupos em situação de vulnerabilidade ; entre outras ;  garantir o acesso dessas comunidades a equipamentos sociais e culturais, bem como a espaços públicos inclusivos e a áreas verdes ;  elaborar planos estratégicos participativos urbanos e ambientais a serem desenvolvidos nas comunidades em situação de vulnerabilidade.
No tocante ao subitem b ), esperava -se que o respondente apresentasse coerência propositiva, englobando os diferentes atores, quais sejam : o governo, em suas diversas instâncias, e a comunidade, também em suas diversas organizações, almejando conter os riscos socioambientais. 5.3.1.3 CONSIDERAÇÕES FINAIS SOBRE A CORREÇÃO DE CONTEÚDO DA QUESTÃO DISCURSIVA DO COMPONENTE DE FORMAÇÃO GERAL De uma maneira geral, os respondentes apresentaram um nível aquém do esperado, em termos de formação geral, destacando -se, neste âmbito, os aspectos a seguir.
Nas respostas relacionadas ao item a ), percebeu -se que, quanto aos riscos, a maioria dos respondentes não se limitou a descrevê -los, ou mesmo a estabelecer um rol deles.
No caso daqueles que descreveram as diversas situações de risco, em sua maioria, as respostas se enquadraram naquilo que estava previsto no padrão de resposta, apresentando os riscos derivados da insegurança relativa ao meio geográfico em que habitam, tais como alagamentos, deslizamentos e doenças por contato com um meio insalubre.
Quanto à descrição do meio, notou -se que houve uma maior aderência por parte dos respondentes, pois a maior parte das respostas, excluindo -se as que foram deixadas em branco, descreveu o meio de alguma forma.
Devido à grande variedade de possíveis descrições do meio geográfico e/ou social em que vive a população brasileira de perfil pobre ou extremamente pobre, o padrão de respostas deixava aberta a possibilidade de considerar respostas variadas em relação às citadas como exemplo em sua redação.
E, de fato, entre as respostas encontradas, verificaram -se muitas que foram consideradas certas, por estarem dentro do perfil previsto como correto, mas que eram distintas dos exemplos específicos apresentados no padrão de respostas.
O mesmo se deu quanto ao item b ) : houve uma grande variedade de proposições para além dos campos temáticos específicos previstos no padrão de respostas, mas que foram consideradas corretas porque atendiam à demanda do item.
Para esse item, o padrão de respostas demandava que as proposições fossem pertinentes ao tema e minimamente desenvolvidas.
No entanto, a maior parte dos respondentes não apresentou proposições desenvolvidas e muitos não fizeram proposições pertinentes ao tema.
O item também demandava que as proposições abrangessem a ação governamental e a participação da comunidade, mas a maior parte dos respondentes envolveu em suas proposições apenas os governos em suas mais diversas instâncias, não abrangendo o papel participativo da comunidade.
Um pequeno grupo de respostas se limitou a apresentar só uma proposição. ao tema.
O item também demandava que as proposições abrangessem a ação governamental e a participação da comunidade, mas a maior parte dos respondentes envolveu em suas proposições apenas os governos em suas mais diversas instâncias, não abrangendo o papel participativo da comunidade.
Um pequeno grupo de respostas se limitou a apresentar só uma proposição.
Dessa forma, o desempenho dos estudantes foi, em geral, muito baixo.
O número de questões deixadas em branco foi considerável, bem como o quantitativo de respostas muito curtas, ocupando menos de cinco das 15 linhas oferecidas para a resposta.
Muitos responderam a apenas um dos dois itens.
A grande maioria dos que buscaram responder ao que a questão demandava, completamente ou parcialmente, não souberam se expressar de modo correto, objetivo e coerente.
Não poucos respondentes repetiram longos trechos do enunciado em suas respostas, demonstrando um baixo nível de formação geral e, também, baixa disposição para a realização da prova.
Muitos respondentes foram desatentos e escreveram a resposta à questão discursiva de componente de Conhecimento Específico no campo dedicado à questão discursiva de Formação Geral, tendo sido incluídos dentre os respondentes que fugiram ao tema.
Uma quantidade proporcionalmente pequena, mas grande numericamente, foi a de estudantes que apresentaram respostas que desrespeitavam os direitos humanos, em sua maioria propondo o controle reprodutivo sobre a população pobre ou extremamente pobre como solução governamental para as situações de risco socioambiental comumente vivenciadas pelas populações brasileiras pobres ou extremamente pobres.
Alguns estudantes também apontaram uma relação necessária entre pobreza e ignorância ou mesmo entre pobreza e violência.
Entre os estudantes que responderam com maior dedicação, buscando dar conta de toda a questão, houve um número significativo que apresentou ter grandes dificuldade de expressão textual, demonstrando que seu baixo domínio da linguagem verbal comprometia a capacidade de produzir respostas adequadas.
Muitas respostas não apresentaram o mínimo de lógica ou não buscaram estabelecer relações entre os conteúdos mobilizados.
Portanto, em sua maioria, os estudantes não apresentaram a capacidade de expressar ideias de maneira clara, coesa e argumentativa por meio da produção de textos, uma habilidade fundamental ao concluinte do ensino Superior.
Outro perfil de destaque foi o de estudantes que pontuaram elementos válidos para responder à questão sem utilizar os conectivos e estabelecer as relações causais.
Também merecem destaques os respondentes que falaram sobre aspectos históricos da pobreza e a descrição dos tipos sociais mais envolvidos em situação de pobreza ou extrema pobreza, sem estabelecer uma relação causal com as situações de risco socioambiental, portanto, não respondendo o que a questão solicitava.
Essas situações demonstraram um perfil majoritário com baixa capacidade de produzir para o item de Formação Geral uma resposta adequada esperada para estudantes concluintes do Ensino Superior.
Os estudantes com perfis de respostas minimamente adequadas ou boas corresponderam ao segundo perfil majoritário.
Em sua maior parte, esse perfil de estudantes respondeu bem ao subitem a ), mas não respondeu bem ou subitem b ), ou vice-versa.
O perfil minoritário foi o de estudantes de nível muito bom a ótimo.
Esses com competência para descrever objetivamente os pontos essenciais da questão, estabelecendo o nexo causal entre eles.
Quanto à relação com o conteúdo, a maior parte dos respondentes soube pensar nos elementos e nas imagens do problema levantado na questão, mas não soube fazer as ligações que estimulam a relação entre esses elementos.
Portanto, demonstraram uma baixa capacidade de inter-relacionar objetos de conhecimento nas diferentes áreas, outra habilidade esperada de um concluinte do Ensino Superior.
Uma grande parte dos estudantes vê o saneamento básico precário como a maior fonte de riscos socioambientais para a população pobre ou extremamente pobre.
As descrições do meio físico, como as situações de moradias em encostas de morros e beiras de rio, também constituíram entre as mais comuns.
Em menor número, mas ainda com relativa frequência, apareceram os casos de estudantes que apontaram o meio como lugar de domínio territorial de criminosos ou os casos que indicaram o meio social como opressivo aos negros.
Quanto às proposições, os estudantes, em sua maioria, não souberam criar estratégias de ação, selecionar métodos e formular propostas.
Uma quantidade considerável de estudantes apresentou proposições muito genéricas como “ construir escolas ”.
A maioria não apontou conexão entre as propostas e a minimização dos riscos socioambientais, como se a ação proposta já evidenciasse o que ela está solucionando, assim também indicando uma grande concentração de estudantes que não aparenta dominar habilidades importantes para concluintes do Ensino Superior. 5.3.1.4 ANÁLISE DE LÍNGUA PORTUGUESA DA QUESTÃO DISCURSIVA DO COMPONENTE DE FORMAÇÃO GERAL Os dados de Nutrição obtidos a partir das respostas à questão discursiva do componente de Formação Geral, no que tange ao desempenho linguístico, são apresentados na Tabela 5.15.
Os dados da Tabela 5.15 mostram que, no quesito desempenho linguístico, os estudantes, de todo o Brasil, obtiveram média de 62,6.
A maior média em relação ao desempenho linguístico foi obtida na região Sudeste ( 63,0 ), e a menor, na região Norte ( 61,5 ).
Quanto à variabilidade das notas, o desvio-padrão de todo o Brasil foi 22,8.
O menor desvio-padrão foi obtido na região Sudeste ( 22,4 ), e o maior desvio-padrão foi obtido na região Norte ( 24,3 ).
A mediana das notas de Língua Portuguesa foi 67,5 para o Brasil, como um todo, a mesma obtida em todas as regiões.
A nota máxima para todo o Brasil foi 100,0, com, pelo menos, um estudante com essa nota em todas as regiões.
A nota mínima foi 0,0 em todas as regiões do país.
No Gráfico 5.8, mostra -se a distribuição das notas de desempenho linguístico do componente de Formação Geral.
O Gráfico 5.8 demonstra que 7,5 \% dos participantes correspondem ao intervalo [ 0 ; 10 ].
A moda dessa distribuição ocorre no intervalo ( 70 ; 80 ], com uma frequência de 25,9 \% dos participantes. 5.3.1.5 COMENTÁRIOS SOBRE A CORREÇÃO DE LÍNGUA PORTUGUESA DAS RESPOSTAS À QUESTÃO DISCURSIVA Os comentários feitos neste item a respeito da correção de Língua Portuguesa das respostas emitidas para a questão discursiva do componente de Formação Geral estão agrupados em três pontos : avaliação da questão, padrão de resposta utilizado na avaliação e análise das respostas dadas pelos estudantes.
A linguagem da questão e do texto motivador foi considerada condizente com o domínio linguístico que se espera de estudantes concluintes de cursos em nível de graduação, sendo -lhes, portanto, acessível.
Do ponto de vista linguístico e textual, o nível da questão foi considerado mediano.
Por meio de um comando claro e objetivo, a questão solicitava que o estudante explicasse determinada relação, exposta no texto motivador, e apresentasse duas propostas para atenuar o problema discutido.
Para compreender o que deveria ser feito, o estudante necessitava ter domínio da leitura e interpretação textual, de modo a entender o enunciado, estabelecer a relação proposta na questão e apontar possíveis ações governamentais e da comunidade para resolver os problemas elencados.
O fato de a proposta textual ser apresentada em dois itens ( a e b ) separadamente levou muitos estudantes a produzirem o texto em dois parágrafos bem definidos e, na maioria dos casos, sem conexão textual entre ambos, ou seja, respondendo à questão em dois ou mais parágrafos que indicavam a letra do item ao qual a resposta se dirigia.
Foram relativamente poucas as provas em que os estudantes respondiam à questão em um texto que articulasse textualmente os dois itens.
Essa formatação dificultou a avaliação do quesito Aspectos Textuais no que se refere à organização textual.
Como a questão demandava certa capacidade de análise, a resposta exigia, em termos textuais, que o estudante fosse além do texto motivador, o que, em parte, passou a identificar a qualidade textual da resposta, ainda que tomasse como base o próprio enunciado da questão e empregasse dados do texto motivador.
Textos com qualidade mediana, em geral, consistiram na retextualização da questão e/ou de enunciados do texto motivador, com alguma autonomia linguística para elaborar argumentação superficial.
A simples cópia – às vezes, com erros ortográficos ou morfossintáticos – da questão ou de enunciados do texto motivador, praticamente sem autonomia textual ou argumentativa, foi um recurso característico de respostas de baixa qualidade em termos linguísticos, porém não penalizado na avaliação do desempenho linguístico.
Em relação ao padrão de resposta utilizado na avaliação, foram considerados os seguintes aspectos : ( 1 ) aspectos ortográficos ; ( 2 ) aspectos textuais ; ( 3 ) aspectos morfossintáticos e de propriedade vocabular adequada à modalidade escrita formal da Língua Portuguesa, exigida pela situação comunicativa.
A seguir, são apresentados cada um dos aspectos avaliados : 1 ) Aspectos ortográficos – referem -se ao domínio das convenções ortográficas da norma-padrão da Língua Portuguesa, como : grafia de vogais e consoantes, uso de maiúsculas e minúsculas, emprego do hífen e da acentuação gráfica, segmentação de palavras. 2 ) Aspectos textuais – referem -se ao domínio dos procedimentos de estruturação textual : organização interna dos períodos, emprego de conectores para a articulação lógica entre os períodos e entre os parágrafos, emprego de marcas de referenciação lexical e pronominal ; e utilização dos sinais de pontuação que contribuem para a organização lógica da frase. 3 ) Aspectos morfossintáticos e de propriedade vocabular – referem -se ao domínio das regras de caráter morfossintático estabelecidas como modelares do ponto de vista da modalidade escrita formal da norma-padrão da Língua Portuguesa : concordância nominal e verbal, regência nominal e verbal, colocação pronominal, flexão nominal e verbal, correlação entre tempos e modos verbais, ausência de marcas de oralidade.
Outra competência combinada a esta na análise das respostas foi a seleção vocabular adequada à modalidade escrita formal da Língua Portuguesa, tendo em vista a interseção entre as duas do ponto de vista das exigências do registro formal da modalidade escrita da norma-padrão.
No que concerne à análise das respostas dadas pelos estudantes, especificamente, em relação ao desempenho linguístico, foram observados muitos deslizes gramaticais no que diz respeito à acentuação e grafia correta das palavras, embora, de maneira geral, os estudantes tenham demonstrado domínio da ortografia vigente.
Quanto aos aspectos textuais, verificou -se predominância de textos organizados sem a presença das três partes esperadas para um texto dissertativo-argumentativo : introdução ; desenvolvimento e conclusão.
A argumentação foi predominantemente rasa e pouco elaborada, havendo frequentes desvios de pontuação, nomeadamente quanto ao emprego da vírgula.
Por fim, quanto aos aspectos morfossintáticos, observou -se um bom domínio das regras de regência e concordância na maior parte dos textos.
A adequação lexical também foi pertinente a uma redação formal na maior parte das vezes.
Em relação ao item a ), destacaram -se como respostas a relação entre comunidades negras periféricas, além de ribeirinhos e indivíduos em vulnerabilidade socioeconômica como sendo o perfil da população brasileira mais propenso a riscos socioambientais.
No que diz respeito ao item b ), foram encontradas com maior frequência propostas relacionadas a atuações governamentais, familiares e sociais diversas, como a necessidade de maiores investimentos e ofertas de trabalho na periferia.
Houve a ocorrência de respostas fracas e incompletas.
O item a ) foi respondido, na maior parte dos textos, de modo discursivo, enquanto o item b ) figurou, por vezes, como uma lista de ações que poderiam resolver o problema.
Em alguns casos, não tão raros, foram produzidos textos nos quais os estudantes se recusaram explicitamente a produzir uma resposta adequada, alegando que a questão era pertinente a uma área de formação diferente da sua.
As respostas encontradas e mesmo a ausência de respostas nos permitem tecer algumas reflexões.
Há constante falta de atenção na leitura do enunciado da questão, o que levou alguns estudantes a responderem à questão específica em local diverso do esperado, prejudicando a resposta.
Por vezes, alguns candidatos produziram as respostas para as duas questões ( Formação Geral e Conhecimento Específico ) no mesmo espaço, prejudicando o desempenho textual e argumentativo.
Seja por desatenção, seja por descuido, a dificuldade de interpretação é sintomática e deve ser considerada.
Desvios ortográficos, morfossintáticos e de pontuação frequentes apontam para lacunas no que concerne à alfabetização e ao desenvolvimento da escrita ainda nos anos do ensino fundamental.
Contudo, o número baixo de textos que tenham obtido notas muito baixas, como 1,0 e 2,0, revela que, entre os estudantes avaliados, não havia pessoas com níveis extremamente baixos de alfabetização e letramento.
Problemas quanto à organização do texto dissertativo-argumentativo, com ausência de introdução, desenvolvimento e conclusão, mesmo que minimamente, apontam para a necessidade de um trabalho mais profundo com a escrita em cursos superiores.
A argumentação rasa ou mesmo o desinteresse por temas sociais revelam o fato de que alguns dos cursos estudados não trabalham o senso crítico do estudante, não o conduzindo a uma vivência plenamente crítica e democrática em sociedade.
A dificuldade em responder ao item b ) da questão demonstrou uma fragilidade teórica na formação, principalmente no que concerne à atuação do profissional em relação a questões de cunho socioeconômico e ao atendimento ao público mais fragilizado socioeconomicamente.
A baixa frequência de desvios morfossintáticos pode ser justificada pela própria competência linguística do estudante, demonstrando que um trabalho mínimo com questões gramaticais nos cursos superiores poderia sanar grande parte do problema.
Relativamente poucos estudantes demonstraram incapacidade quanto à escolha de palavras adequadas para a produção de um texto formal, apresentando competência genérica, na maior parte das vezes, satisfatória.
Textos que obtiveram nota máxima foram raros, o que aponta para a necessidade de reforçar a importância do trabalho com a leitura e escrita em âmbito formal com os estudantes dos cursos avaliados. 5.3.1.6 CONSIDERAÇÕES FINAIS SOBRE A CORREÇÃO DE LÍNGUA PORTUGUESA DAS RESPOSTAS À QUESTÃO DISCURSIVA DO COMPONENTE DE FORMAÇÃO GERAL A análise global em relação ao desempenho linguístico dos cursos avaliados no Enade 2023 revelou um desempenho mediano dos estudantes avaliados.
Quanto aos aspectos textuais e argumentativos, verificou -se um nível baixo de conhecimento e mesmo de interesse, por parte de muitos estudantes, em refletir sobre os problemas sociais envolvendo pessoas em situação de vulnerabilidade e fenômenos de risco socioambiental.
A maioria dos estudantes não conseguiu apresentar uma análise crítica, tampouco responder conforme solicitado pela questão.
Alguns estudantes não foram atentos às informações apresentadas nos textos nem ao que os enunciados solicitavam.
Alguns demonstraram o domínio sobre o assunto ou, no mínimo, o comprometimento e o interesse para responder as questões.
Entende -se que o tempo, a quantidade de questões e a pressão podem influenciar no momento de elaborar a resposta.
Além disso, alguns estudantes de cursos específicos não se sentiram motivados a falar de um assunto \textbf{julgado} como “ social ” e não pertinente à sua área de atuação profissional, o que é um forte representativo de vieses ideológicos que perpassam determinadas formações.
Cabe ressaltar que a questão abordava um conteúdo atual e relevante a ser considerado por um indivíduo plenamente capaz de agir em sociedade de forma crítica, de qualquer área de formação em nível superior.
Destaca -se que estudantes com aparente domínio da área de saúde apresentaram inclinação crítica a abordar o tema de forma bastante humanizada.
Quanto aos aspectos ortográficos e gramaticais de natureza morfossintática, observam -se muitas lacunas trazidas desde a educação básica.
Entende -se que a universidade não pode sanar todos os problemas de cunho ortográfico e gramatical, mas é preciso salientar a importância de disciplinas voltadas à leitura e produção de texto serem incluídas como disciplinas obrigatórias nos cursos avaliados. 5.3.2 COMPONENTE DE CONHECIMENTO ESPECÍFICO As análises dos resultados de desempenho dos estudantes de Nutrição, na questão discursiva relativa ao componente de Conhecimento Específico, encontram -se na Tabela 5.16 e no Gráfico 5.9. 5.3.2.1 ANÁLISE DA QUESTÃO DISCURSIVA DO COMPONENTE DE CONHECIMENTO ESPECÍFICO A Tabela 5.16 apresenta dados sobre a questão discursiva do componente de Conhecimento Específico.
A partir dos dados registrados na Tabela 5.16, observa -se que a média foi mais baixa se comparada à média obtida na questão discursiva do componente de Formação Geral.
Enquanto, no componente de Formação Geral, a média para estudantes de Nutrição de todo o Brasil foi 42,4 ; na parte de Conhecimento Específico a média foi 27,5.
A maior média desse componente foi obtida pelos estudantes da região Centro-Oeste ( 29,4 ), e a menor, pelos da região Sul ( 26,8 ).
Quanto à variabilidade das notas, o desvio-padrão de todo o Brasil foi 24,7.
O maior desvio-padrão foi encontrado na região Nordeste ( 25,1 ), e o menor nas regiões Norte e Sudeste ( 24,5 ).
A nota máxima obtida foi 100,0 por pelo menos um estudante em todas as regiões.
A nota mínima foi 0,0 em todas as regiões do país.
No Gráfico 5.9, apresenta -se a distribuição das notas na questão discursiva no componente de Conhecimento Específico.
À luz dos dados do Gráfico 5.9, verifica -se que a moda da distribuição das notas da questão discursiva do componente de Conhecimento Específico ocorre no intervalo [ 0 ; 10 ], com 31,4 \% do total de participantes.
Observa -se um valor considerável no intervalo [ 0 ; 10 ], cerca de 31,4 \% dos participantes, entretanto, vale ressaltar que, nesse intervalo, são contabilizadas também as questões deixadas em branco. 5.3.2.2 COMENTÁRIOS SOBRE AS RESPOSTAS À QUESTÃO DISCURSIVA A questão discursiva apresentada para área de Nutrição foi composta por dois textos motivadores e três itens ( a, b e c ).
Foi uma questão de dificuldade mediana, em que o primeiro texto motivador exigia um resgate e uma reflexão sobre a importância da Educação Alimentar e Nutricional ( EAN ) para as políticas públicas, em que o foco devia ser a valorização do alimento em sua totalidade, considerando os aspectos sociais, culturais, simbólicos, psicológicos, afetivos, sensoriais, econômicos, geográficos e ambientais.
O foco não devia ser o nutriente, mas, sim, os alimentos, os saberes e as diferentes formas de preparo, destacando o desenvolvimento das habilidades culinárias como prática para que os indivíduos tenham autonomia e possam fazer escolhas alimentares saudáveis, de acordo com o ambiente em que estão inseridos e com suas histórias.
Em outras palavras, a EAN deve abordar essas múltiplas dimensões em suas ações \textbf{formativas}, de prática contínua e permanente, transdisciplinar, multifatorial e intersetorial, que promovam modos de vida saudáveis, na perspectiva da segurança alimentar e nutricional.
O texto motivador 2 apresentou um artigo do Fundo Nacional de Desenvolvimento da Educação ( FNDE ) que falava sobre o atendimento da alimentação escolar aos estudantes da educação básica no âmbito do Programa Nacional de Alimentação Escolar ( PNAE ).
O artigo contextualizava que, para a promoção da alimentação saudável nas escolas e o fortalecimento da agricultura familiar e dos empreendimentos familiares rurais e/ou suas organizações, o planejamento adequado de cardápio pelo nutricionista Responsável Técnico ( RT ) vinculado à entidade executora deve respeitar os hábitos alimentares e os aspectos culturais locais ; que os alimentos que o compõe devem ser, em sua maioria, in natura e minimamente processados, considerando a diversificação agrícola da região nas diferentes estações do ano e serem provenientes da agricultura familiar e/ou dos empreendimentos rurais locais e/ou das suas associações.
Para isso, espera -se que o RT adote ações para o diálogo com as organizações locais da agricultura familiar e com as entidades de assistência técnica e extensão rural do município/estado e, ainda, que incentive a participação dos agricultores familiares para o fornecimento de alimentos para a alimentação escolar, como grupo formal, informal ou fornecedor individual.
Destaca -se, também, como função profissional de nutrição nas escolas a prescrição dietética e a adaptação do cardápio para atender o estudante com condição clínica específica, de forma que ele se sinta incluído no ambiente escolar, contribuindo para a promoção e a manutenção da saúde e qualidade de vida.
No caso em questão, havia a necessidade de orientação nutricional individualizada para as crianças com DM1 ( Diabetes Mellitus tipo 1 ), com avaliação nutricional e prescrição dietética de acordo com a dose de insulina administrada, além de orientações específicas e atividades de educação alimentar e nutricional para a equipe escolar, estudantes e pais.
Considerando a abrangência da questão, a qual poderia ser respondida de múltiplas formas, as respostas foram diversas, além de muitas vezes direcionadas somente para a população de baixa renda.
Alguns estudantes conseguiram responder de forma parcial ou total as questões, porém muitos desviaram completamente do foco, zerando todos os itens.
Muitos estudantes limitaram suas explicações, no item a ), ao aumento do consumo de alimentos in natura, integrais e fibras e à redução do consumo de alimentos ultraprocessados, demonstrando que desconheciam os princípios do Marco de Referência de Educação Alimentar e Nutricional.
Observou -se o mesmo padrão em relação ao PNAE, sendo que o desconhecimento das diretrizes do programa prejudicou as respostas do item b ).
Nos enunciados dos itens b ) e c ) foram solicitadas duas ações, mas muitos respondentes citaram apenas uma ação em cada item.
Foi possível notar que os estudantes apresentaram dificuldade na interpretação da questão e na organização da resposta.
Em muitos casos, não indicaram o que era a resposta ao item a ), b ) e c ), escrevendo um único texto que supostamente responderia aos três itens.
Todavia, os avaliadores consideraram as respostas e atribuíram os conceitos equivalentes a cada item.
No item b ), muitos abordaram, de forma equivocada, o incentivo à criação de hortas escolares e comunitárias, cultivadas pelas próprias famílias ; a aquisição de alimentos in natura pelo próprio nutricionista em feiras locais ; a realização de palestras com os responsáveis sobre a importância da alimentação saudável e da agricultura familiar ; a fiscalização do processo de compra e ações educativas para incentivo do aproveitamento integral dos alimentos no ambiente escolar.
No item c ), observou -se muitas respostas sobre a dietoterapia específica para quem tem diabetes sem levar em consideração a assistência nutricional do escolar no âmbito do PNAE.
Também foi recorrente a sugestão de palestras de alimentação saudável e cardápios saudáveis ( como, por exemplo, retirando farinha branca e carboidratos simples ), compra de alimentos diet para a alimentação escolar, encaminhamento para o médico e incentivo à prática de atividade física.
Diante do exposto, a percepção é a de que muitos estudantes desconhecem a aplicação das diretrizes do Programa Nacional de Alimentação Escolar ( PNAE ) na prática.
A percepção, a partir de uma análise qualitativa, é de que o desempenho dos estudantes foi ruim.
Muitos textos contemplaram apenas parcialmente a questão, respondendo apenas ao item c ), mencionando a “ adaptação do cardápio ”.
A ausência ou dificuldade em atender ao que foi solicitado nos itens a ) e b ) da questão pode evidenciar uma fragilidade na compreensão crítica do conteúdo, uma vez que muitos estudantes não abordaram ou discorreram de forma inadequada sobre os princípios do Marco de Referência em EAN e/ou sobre as diretrizes do PNAE.
Além disso, foi identificado também uma dificuldade dos estudantes em estabelecer estratégias do tratamento dietoterápico do diabetes. 5.3.2.3 CONSIDERAÇÕES FINAIS SOBRE A QUESTÃO DISCURSIVA DO COMPONENTE DE CONHECIMENTO ESPECÍFICO Em relação à questão discursiva da área de Nutrição, a análise global revelou um nível entre baixo e mediano de conhecimento.
Poucos estudantes demonstraram o domínio sobre o assunto.
Dessa forma, ficou evidente o uso de termos do jargão profissional e do senso comum em detrimento dos termos técnicos da área, bem como uma dificuldade de interpretação de texto e de escrita legível por alguns respondentes.
Destaca -se, ainda, que foi recorrente a presença de erros ortográficos e gramaticais nos textos avaliados.
Alguns estudantes utilizaram o espaço reservado à resposta da questão para desabafos e reclamações quanto à instituição de ensino da qual fazem parte.
As respostas encontradas, e mesmo a ausência de respostas, nos permitem fazer duas reflexões : a primeira refere -se aos erros de interpretação, que foram uma marca das respostas a essa questão.
Parte das respostas abordou as recomendações do Guia Alimentar para a População Brasileira e ações de EAN de forma abrangente.
Esses erros de interpretação do enunciado podem indicar uma fragilidade da formação.
Seja por desatenção ou descuido, a dificuldade de interpretação é sintomática e deve ser considerada.
A segunda é em relação ao fato de os estudantes não contextualizarem as ações do nutricionista, no âmbito do PNAE, para a aquisição de gêneros alimentícios da agricultura familiar e na assistência nutricional no DM1.
A dificuldade em responder aos itens b ) e c ) demonstra uma fragilidade teórica na formação relacionada ao PNAE e em conceitos de outras áreas, principalmente a relação entre a Nutrição Clínica e a Saúde Coletiva.

% matched lemmas: formativo, julgar
\end{textsample}
